\section{Additional Theory Details and Proofs}
\label{app:theory-proofs}

This appendix supplies the proof details that are not central to the main
paper narrative. All results concern oracle-start local defects for
deterministic ODE solvers and support local risk calibration only.

\subsection{Mixed residual decomposition}
\label{app:mixed-residual}

Let $T=T_{k,G}(b)$ satisfy $T^\top G(b)T=1$. The $G(b)$-orthogonal tangent
projection of $r$ is
\begin{equation}
  r_\parallel^{G(b)}=(T^\top G(b)r)T,
\end{equation}
and $r_\perp^{G(b)}=r-r_\parallel^{G(b)}$. Orthogonality gives
\begin{equation}
  \|r\|_{G(b)}^2
  =
  \|r_\perp^{G(b)}\|_{G(b)}^2
  +
  \|r_\parallel^{G(b)}\|_{G(b)}^2,
\end{equation}
with
\begin{equation}
  \|r_\parallel^{G(b)}\|_{G(b)}^2=(T^\top G(b)r)^2.
\end{equation}
Substituting into Eq.~\eqref{eq:mixed-norm} yields
Eq.~\eqref{eq:mixed-decomposition}.

\subsection{Equal monitor mass}
\label{app:equal-monitor-mass}

For completeness we prove Lemma~\ref{lem:equal-monitor-mass}. The interval
masses satisfy $\sum_{i=0}^{N-1}M_i=\Omega_{\mathcal R}$, so their maximum is
at least their average. Equality holds exactly when all masses are equal. Since
$\omega_{\mathcal R}>0$, the cumulative monitor clock
$\tau_{\mathcal R}$ is strictly increasing and invertible, and the inverse-CDF
construction in Eq.~\eqref{eq:inverse-cdf-schedule} gives equal increments of
normalized cumulative mass.

\subsection{Full proof of asymptotic minimax optimality}
\label{app:minimax-proof}

For the equal-mass schedule $U^\Omega$,
\begin{equation}
  \mathcal M_{\mathcal R}([u_i^\Omega,u_{i+1}^\Omega])
  =
  \frac{\Omega_{\mathcal R}}{N}.
\end{equation}
Since $\omega_{\mathcal R}(u)\ge\omega_{\min}$,
\begin{equation}
  h_i^\Omega
  =
  u_{i+1}^\Omega-u_i^\Omega
  \le
  \frac{\Omega_{\mathcal R}/N}{\omega_{\min}}
  =
  O(N^{-1}).
\end{equation}
Using Assumption~\ref{assump:uniform-monitor},
\begin{align}
  \bar D_{\mathcal R}(u_i^\Omega,u_{i+1}^\Omega)
  &\le
  \mathcal M_{\mathcal R}([u_i^\Omega,u_{i+1}^\Omega])^q
  +C_E(h_i^\Omega)^{q+1}\\
  &\le
  \left(\frac{\Omega_{\mathcal R}}{N}\right)^q
  +C_1N^{-(q+1)}.
\end{align}
Taking the maximum over intervals proves the upper bound.

Now consider any admissible schedule $U$ with $h_{\max}\le C_h/N$. Define
$M_i=\mathcal M_{\mathcal R}([u_i,u_{i+1}])$. Since
$\sum_iM_i=\Omega_{\mathcal R}$, there exists $j$ such that
$M_j\ge\Omega_{\mathcal R}/N$. Again by
Assumption~\ref{assump:uniform-monitor},
\begin{align}
  \bar D_{\mathcal R}(u_j,u_{j+1})
  &\ge
  M_j^q-C_Eh_j^{q+1}\\
  &\ge
  \left(\frac{\Omega_{\mathcal R}}{N}\right)^q
  -C_E\left(\frac{C_h}{N}\right)^{q+1}\\
  &=
  \left(\frac{\Omega_{\mathcal R}}{N}\right)^q
  -C_2N^{-(q+1)}.
\end{align}
Therefore the maximum interval defect of any admissible schedule is at least
this lower bound.

\subsection{Monte Carlo estimation of the mean monitor}
\label{app:mc-proof}

Since $x_k^\star(u)\sim p_u$, the random variables
$a_k(u)=a_{\mathcal S}(x_k^\star(u),u)$ are i.i.d. samples from the
distribution of $a_{\mathcal S}(x,u)$ under $p_u$. If the variance is finite,
the central limit theorem gives
\begin{equation}
  \frac{1}{K}\sum_{k=1}^K a_k(u)
  -
  \mathbb E_{x\sim p_u}[a_{\mathcal S}(x,u)]
  =
  O_p(K^{-1/2}).
\end{equation}
If coefficients are estimated from finite probes, write
$\widehat a_k(u)=a_k(u)+\delta_k(u)$ with
$K^{-1}\sum_k\delta_k(u)=O(\varepsilon_{\rm probe}(u))$. Then
\begin{equation}
  \widehat a_{\rm mean}(u)-\bar a_{\rm mean}(u)
  =
  O_p(K^{-1/2})+O(\varepsilon_{\rm probe}(u)).
\end{equation}

\subsection{Inverse-CDF stability}
\label{app:inverse-stability-proof}

Let
\begin{equation}
  F(u)=\int_{u_{\min}}^u\omega_{\mathcal R}(\xi)d\xi,
  \qquad
  \widehat F(u)=\int_{u_{\min}}^u\widehat\omega_{\mathcal R}(\xi)d\xi.
\end{equation}
If $\sup_u|\widehat F(u)-F(u)|\le\varepsilon_F$ and
$\omega_{\mathcal R}(u)\ge\omega_{\min}>0$, then $F^{-1}$ is Lipschitz with
constant at most $1/\omega_{\min}$. The corresponding exact and estimated
schedules satisfy
\begin{equation}
  \max_m|\widehat u_m-u_m|
  \le
  \frac{\varepsilon_F}{\omega_{\min}}+O(\varepsilon_F^2).
\end{equation}
The second-order term accounts for local nonlinear curvature of the inverse
map and normalization effects in the estimated cumulative mass.

\subsection{Distributional coupling interpretation}
\label{app:w2-coupling}

Let $p_{\star,T}$ be the terminal teacher distribution and
$p_{\mathcal S,T}$ the terminal student distribution induced by schedule $U$.
For a shared initial sample $x_0$, define
\begin{equation}
  x_T^\star=F_{u_{\min}\rightarrow T}(x_0),
  \qquad
  x_T^{\mathcal S}=\Phi_{\mathcal S,U}(x_0).
\end{equation}
This shared-noise construction is a coupling between $p_{\star,T}$ and
$p_{\mathcal S,T}$, hence
\begin{equation}
  W_2^2(p_{\mathcal S,T},p_{\star,T})
  \le
  \mathbb E_{x_0}\!\left[\|x_T^{\mathcal S}-x_T^\star\|_{G_T}^2\right].
\end{equation}
Under a first-order error propagation expansion,
\begin{equation}
  x_T^{\mathcal S}-x_T^\star
  =
  \sum_{i=0}^{N-1}S_ie_i+\text{higher-order terms},
\end{equation}
where $e_i$ is the oracle-start local error at step $i$, stability of the
downstream flow gives the conservative proxy
\begin{equation}
  W_2^2(p_{\mathcal S,T},p_{\star,T})
  \le
  C_{\rm stab}(U)
  \sum_{i=0}^{N-1}
  \mathbb E_{x\sim p_{u_i}}
  [d_{\mathcal S}(x;u_i,u_{i+1})]
  +\text{higher-order terms}.
\end{equation}
This is an interpretation, not the main guarantee: local defect controls a
first-order coupling proxy under stability assumptions, but it is not an exact
final image-quality objective.
